\chapter{Derivando ecuaciones diferenciales}
\noindent
Motivando el estudio que sigue, si pensamos nuevamente en el problema del péndulo:
\begin{equation*}
    \ddot{\theta} + \frac{g}{l}\sen\theta = 0, \qquad \left\{\begin{array}{l}
        \theta(0) = \theta_0 \\
        \dot{\theta}(0) = \omega_0
    \end{array}\right.
\end{equation*}
La única solución explícita que conocemos es $\theta(t;0,0) = 0\quad t\in \mathbb{R}$. Para $\theta_0,\omega_0$ pequeños podemos pensar en aproximar $\theta(t;\theta_0,\omega_0)$ a partir de $\theta(t;0,0)$.\\
\ \\

\noindent
Así, el marco teórico en el que trabajamos es que queremos resolver un problema de valores iniciales:
\begin{equation*}
    \dot{x} = X(t,x), \qquad x(t_0) = \xi
\end{equation*}
y tenemos $\varphi(t)$ una solución para la condición $\varphi(t_0) = \xi_0$. Con el fin de resolver el problema anterior a partir de $\varphi(t)$, queremos realizar la aproximación:
\begin{equation*}
    x(t;t_0,\xi) = \underbrace{x(t;t_0,\xi_0)}_{\varphi(t)} + \underbrace{\frac{\partial x}{\partial \xi}(t;t_0,\xi_0)}_{(\ast)}(\xi-\xi_0) + o(\|\xi-\xi_0\|) \qquad (\xi\to\xi_0)
\end{equation*}
La dificultad de esta lección es tratar de entender el término $(\ast)$. Dejando clara la fórmula anterior, debemos ya saber que:
\begin{itemize}
    \item $x(t;t_0,\xi)\in \mathbb{R}^d$.
    \item $x(t;t_0,\xi_0)\in \mathbb{R}^d$.
    \item $\frac{\partial x}{\partial \xi}(t;t_0,\xi_0)\in \mathbb{R}^{d\times d}$.
    \item $(\xi-\xi_0)\in \mathbb{R}^d$
\end{itemize}

\noindent
El siguiente Teorema es llamado en algunos libros ``Teorema de diferenciabilidad de Peano''. Sin embargo, el matemático que lo enunció y demostró fue Poincaré en su tesis doctoral, que lo hizo para ecuaciones analíticas y luego se particularizó para clase 1.

% // TODO: Fijar bien las hipótesis del Teorema
\begin{teo}[de diferenciabilidad respecto a condiciones iniciales]\label{teo:dif_cond}~\\
    \noindent
    Tenemos la ecuación diferencial:
    \begin{equation*}
        \dot{x} = X(t,x)
    \end{equation*}
    donde $X:D\subset\mathbb{R}\times\mathbb{R}^d\to \mathbb{R}^d$ es un campo continuo con $D$ un abierto conexo y de manera\footnote{Observemos que es la misma hipótesis que exigíamos en el Teorema de unicidad de soluciones.} que:
    \begin{equation*}
        \exists \frac{\partial X}{\partial x}(t,x) \qquad \forall (t,x)\in D
    \end{equation*}
    y además la aplicación $(t,x)\in D\longmapsto \frac{\partial X}{\partial x}(t,x)\in \mathbb{R}^{d\times d}$ es continua.\\

    \noindent
    Como tenemos existencia y unicidad de soluciones, podemos considerar la solución general de la ecuación:
    \Func{x}{\cc{D}}{\bb{R}^d}{(t;t_0,x_0)}{x(t;t_0,x_0)}
    de la que sabemos que $\cc{D}$ es abierto y $x:\cc{D}\to\mathbb{R}^d$ es continua.\\

    \noindent
    Bajo estas hipótesis, se tiene que $x\in C^1(\cc{D},\mathbb{R}^d)$.
\end{teo}

\noindent
El Teorema además nos da la forma de calcular las derivadas de $x:\cc{D}\to \mathbb{R}^d$. Antes de dar su fórmula, motivémosla realizando previamente un cálculo formal, que tomará rigurosidad en la prueba del Teorema:
\begin{itemize}
    \item En primer lugar la derivada respecto a $t$:
        \begin{equation*}
            \mathbb{R}^d\ni \frac{\partial x}{\partial t}(t;t_0,x_0) = \dot{x}(t;t_0,x_0) = X(t,x(t;t_0,x_0))
        \end{equation*}
    \item Ahora, la parcial respecto $x_0$:
        \begin{equation*}
            \frac{\partial x}{\partial x_0}(t;t_0,x_0) \in \mathbb{R}^{d\times d}
        \end{equation*}
        Para calcularla, pensamos en la identidad anterior:
        \begin{equation*}
            \frac{\partial x}{\partial t}(t;t_0,x_0) = X(t,x(t;t_0,x_0))
        \end{equation*}
        Y podemos derivar respecto $x_0$:
        \begin{equation*}
            \frac{\partial }{\partial x_0}\left(\frac{\partial x}{\partial t}(t;t_0,x_0)\right) = \frac{\partial X}{\partial x}(t,x(t;t_0,x_0))\frac{\partial x}{\partial x_0}(t;t_0,x_0)
        \end{equation*}
        Ahora, ``mágicamente'' escribimos:
        \begin{equation*}
            \frac{\partial }{\partial x_0}\left(\frac{\partial x}{\partial t}(t;t_0,x_0)\right)  = \frac{\partial }{\partial t}\left(\frac{\partial x}{\partial x_0}(t;t_0,x_0)\right)
        \end{equation*}
        Y si pensamos que $\frac{\partial x}{\partial x_0}(t;t_0,x_0) = Y(t)$ es una incógnita en función del tiempo $t$, tenemos la ecuación:
        \begin{equation*}
            \dot{Y}(t) = \frac{\partial X}{\partial x}(t,x(t;t_0,x_0))\ Y(t)
        \end{equation*}
        Es decir, la matriz $\frac{\partial x}{\partial x_0}(t;t_0,x_0)$ es solución de la ecuación lineal homogénea superior.

        Si pensamos ahora que (tenemos que buscar condición inicial para la ecuación superior) tenemos $x(t_0;t_0,x_0) = x_0$, derivando respecto $x_0$ tenemos que\footnote{Recordemos que $x_0\in \mathbb{R}^d$, con lo que $\frac{\partial x_0}{\partial x_0}=Id_d$.}: 
        \begin{equation*}
            \frac{\partial x}{\partial x_0}(t_0;t_0,x_0) = Id_d
        \end{equation*}
        Así, tenemos la condición $Y(t_0) = Id_d$, por lo que $Y$ es la matriz principal fundamental\footnote{Conviene repasar Ecuaciones Diferenciales I.} en $t_0$ de la ecuación matricial anterior.
    \item Para la parcial respecto $t_0$:
        \begin{equation*}
            \frac{\partial x}{\partial t_0}(t;t_0,x_0) \in \mathbb{R}^d
        \end{equation*}
        Haciendo una cuenta parecida a la anterior:
        \begin{equation*}
            X(t,x) = \frac{\partial x}{\partial t} = \frac{\partial }{\partial t}\left(\frac{\partial x}{\partial t_0}\right) = \frac{\partial X}{\partial x}(t,x)\frac{\partial x}{\partial t_0}
        \end{equation*}
        Y nos sale la misma ecuación que antes, pero ahora $y(t) = \frac{\partial x}{\partial t_0}(t;t_0,x_0)$ es un vector de $\mathbb{R}^d$, solución de la ecuación:
        \begin{equation*}
            \dot{y} = \frac{\partial X}{\partial x}(t,x(t;t_0,x_0))\ y
        \end{equation*}
        Para la que tenemos que buscar una condición inicial, que se obtiene derivando la condición inicial. Tenemos:
        \begin{equation*}
            x(t_0;t_0,x_0)=x_0
        \end{equation*}
        Y queremos derivarlo respecto $t_0$:
        \begin{equation*}
            \frac{\partial x}{\partial t}(t_0;t_0,x_0) + \frac{\partial x}{\partial t_0}(t_0;t_0,x_0) = 0
        \end{equation*}
        Y anteriormente teníamos que $\frac{\partial x}{\partial t}(t;t_0,x_0) = X(t,x(t;t_0,x_0))$, con lo que la condición inicial es:
        \begin{equation*}
            \underbrace{\frac{\partial x}{\partial t}(t_0;t_0,x_0)}_{X(t_0,x_0)} + \underbrace{\frac{\partial x}{\partial t_0}(t_0;t_0,x_0)}_{y(t_0)} = 0 \quad\Longrightarrow\quad
            y(t_0) = -X(t_0,x_0)
        \end{equation*}
\end{itemize}

\noindent
Así, la tesis del Teorema~\ref{teo:dif_cond} anterior es:

\noindent
$x\in C^{1}(\cc{D},\mathbb{R}^d)$ y:
\begin{align*}
    \frac{\partial x}{\partial t}(t;t_0,x_0) &= X(t,x(t;t_0,x_0)) \\
    \frac{\partial x}{\partial x_0}(t;t_0,x_0) &= Y(t), \qquad \text{donde $Y(t)$ es solución matricial de:} \\
                                               &\qquad \qquad\qquad (\ast)\qquad   \dot{Y}(t) = \frac{\partial X}{\partial x}(t,x(t;t_0,x_0))\ Y(t), \quad Y(t_0)  = Id_d \\
    \frac{\partial x}{\partial t_0}(t;t_0,x_0) &= y(t), \qquad \text{donde $y(t)$ es solución de:} \\
                                               &\qquad \qquad\qquad \quad \qquad   \dot{y}(t) = \frac{\partial X}{\partial x}(t,x(t;t_0,x_0))\ y(t), \quad y(t_0)  = -X(t_0,x_0)
\end{align*} 
donde $(\ast)$ es la ecuación variacional o linealizada de la solución $x(t;t_0,x_0)$. 
